\documentclass[11pt,a4paper]{article}
\usepackage{fontspec,amsmath,amssymb,amsthm,microtype}
\setmainfont{lmroman10-regular.otf}[BoldFont=lmroman10-bold.otf,ItalicFont=lmroman10-italic.otf,BoldItalicFont=lmroman10-bolditalic.otf]
\usepackage{xcolor}
\usepackage[most]{tcolorbox}
\usepackage{mathrsfs}
\providecommand{\lt}{<}
\providecommand{\gt}{>}
\setmonofont{DejaVuSansMono.ttf}[BoldFont=DejaVuSansMono-Bold.ttf,ItalicFont=DejaVuSansMono-Oblique.ttf,BoldItalicFont=DejaVuSansMono-BoldOblique.ttf,Scale=0.86]
\newfontfamily\leanfallback{STIXTwoMath-Regular.otf}[Scale=0.86]
\definecolor{leanrule}{gray}{0.62}
\definecolor{leanback}{gray}{0.962}
\definecolor{leantitle}{gray}{0.22}
\definecolor{leanmuted}{gray}{0.38}
\newtcolorbox{leanbox}[1]{enhanced,breakable,colback=leanback,colframe=leanrule,
  boxrule=0.4pt,arc=0.6mm,left=2.4mm,right=2.4mm,top=1.8mm,bottom=1.8mm,
  before skip=0.8em,after skip=1em,fonttitle=\small\bfseries,coltitle=leantitle,
  colbacktitle=leanback,titlerule=0.3pt,title={#1}}
\newlength{\leancharwidth}
\newenvironment{leancode}{\par\smallskip\begingroup\ttfamily\footnotesize\raggedright
  \settowidth{\leancharwidth}{M}\setlength{\parskip}{0pt}\setlength{\parindent}{0pt}
  \frenchspacing}{\par\endgroup\smallskip}
\newcommand{\leanline}[2]{\hangindent=\dimexpr#1\leancharwidth+4\leancharwidth\relax
  \hangafter=1\noindent\hspace*{#1\leancharwidth}#2\par}
\newcommand{\leansource}[1]{\par{\footnotesize\color{leanmuted}#1}\par}
\newcommand{\leanlabel}[1]{\par\smallskip\noindent{\small\bfseries #1}\par}
\usepackage[margin=25mm]{geometry}
\usepackage[unicode,hidelinks]{hyperref}
\hypersetup{pdfauthor={LangevinLean Project}}
\setlength{\emergencystretch}{3em}
\setlength{\parskip}{0.4em}
\setlength{\parindent}{0pt}
\allowdisplaybreaks[2]
\raggedbottom
\title{\texorpdfstring{Moreau envelope: gradient, smoothness and uniform approximation}{Moreau envelope: gradient, smoothness and uniform approximation}}
\author{LangevinLean Project}
\date{}
\begin{document}\maketitle
Draft mathematical exposition; these notes may contain mistakes.


Let \(E\) be a finite-dimensional real inner-product space, \(f:E\to\mathbb R\) convex and bounded below, \(\lambda>0\), and \(P_\lambda\) the proximal map of the previous section. The \emph{Moreau envelope}


\[f^\lambda(x)=\min_{p\in E}\Bigl(f(p)+\frac{\|p-x\|^2}{2\lambda}\Bigr)=f(P_\lambda(x))+\frac{\|P_\lambda(x)-x\|^2}{2\lambda}\]


is shown to be convex and continuously differentiable with gradient \(\nabla f^\lambda(x)=r_\lambda(x)=(x-P_\lambda(x))/\lambda\), and this gradient is \(\lambda^{-1}\)-Lipschitz. The proof rests on the two-sided bound \(0\le f^\lambda(y)-f^\lambda(x)-\langle r_\lambda(x),y-x\rangle\le\|y-x\|^2/(2\lambda)\); no differentiability of \(f\) is assumed.




If moreover \(f\) is \(K\)-Lipschitz, then \(\|\nabla f^\lambda\|\le K\), \(\|x-P_\lambda(x)\|\le\lambda K\) and \(0\le f-f^\lambda\le\lambda K^2/2\) on all of \(E\). The gradient formula identifies the smoothed particle algorithms of the later sections with Euler schemes for \(g_1+g_2^\lambda\), and the uniform gap \(\lambda K^2/2\) controls the distance between the smoothed and the original marginal maximizers.


\section{\texorpdfstring{The envelope, the residual and two value bounds}{The envelope, the residual and two value bounds}}


\textbf{Setting.} In this section \(E\) is a finite-dimensional real inner-product space, \(f:E\to\mathbb R\) is convex and bounded below, and \(\lambda>0\). By the previous section every \(x\in E\) has a unique proximal point \(P_\lambda(x)\), the minimizer of \(Q_{\lambda,x}(p)=f(p)+\|p-x\|^2/(2\lambda)\). The \emph{Moreau envelope} of \(f\) with parameter \(\lambda\) and the \emph{proximal residual} are


\[f^\lambda(x)=Q_{\lambda,x}\bigl(P_\lambda(x)\bigr)=\min_{p\in E}Q_{\lambda,x}(p),\qquad r_\lambda(x)=\frac{x-P_\lambda(x)}{\lambda}.\]


\textbf{Claim.} (a) \(f^\lambda(x)\le f(x)\) for every \(x\in E\). (b) If \(b\in\mathbb R\) satisfies \(f(z)\ge b\) for every \(z\in E\), then \(Q_{\lambda,x}(p)\ge b\) for all \(x,p\in E\); in particular \(f^\lambda(x)\ge b\). Part (b) uses only \(\lambda>0\) and the lower bound.




\textbf{Proof.} (a) \(P_\lambda(x)\) minimizes \(Q_{\lambda,x}\); the competitor \(p=x\) gives \(f^\lambda(x)\le Q_{\lambda,x}(x)=f(x)+0=f(x)\). (b) Since \(\|p-x\|^2\ge0\) and \(2\lambda>0\), \(Q_{\lambda,x}(p)=f(p)+\|p-x\|^2/(2\lambda)\ge f(p)\ge b\). Taking \(p=P_\lambda(x)\) gives \(f^\lambda(x)\ge b\). \(\square\)




Thus \(f^\lambda\) is a real-valued function with \(b\le f^\lambda\le f\).

\begin{leanbox}{Lean: \texttt{moreau}, \texttt{residual}, \texttt{moreau\_\allowbreak{}le} and 1 more}
\leansource{\texttt{moreau} \quad ProximalLangevin/\allowbreak{}Moreau.lean, lines 24–25}
\begin{leancode}
\leanline{0}{def\ moreau\ (f\ :\ E\ →\ ℝ)\ (lam\ :\ ℝ)\ (x\ :\ E)\ :\ ℝ\ :=}
\leanline{2}{proxObjective\ f\ lam\ x\ (prox\ f\ lam\ x)}
\end{leancode}
\leansource{\texttt{residual} \quad ProximalLangevin/\allowbreak{}Moreau.lean, lines 28–29}
\begin{leancode}
\leanline{0}{def\ residual\ (f\ :\ E\ →\ ℝ)\ (lam\ :\ ℝ)\ (x\ :\ E)\ :\ E\ :=}
\leanline{2}{lam⁻¹\ •\ (x\ -\ prox\ f\ lam\ x)}
\end{leancode}
\leansource{\texttt{moreau\_\allowbreak{}le} \quad ProximalLangevin/\allowbreak{}Moreau.lean, lines 37–38}
\begin{leancode}
\leanline{0}{theorem\ moreau\_le\ (x\ :\ E)\ :\ moreau\ f\ lam\ x\ ≤\ f\ x}
\end{leancode}
\leansource{\texttt{le\_\allowbreak{}moreau} \quad ProximalLangevin/\allowbreak{}Moreau.lean, lines 42–44}
\begin{leancode}
\leanline{0}{theorem\ le\_moreau\ \{b\ :\ ℝ\}\ (hbf\ :\ ∀\ z,\ b\ ≤\ f\ z)\ (x\ :\ E)\ :\ b\ ≤\ moreau\ f\ lam\ x}
\end{leancode}
\leanlabel{How the Lean proof follows the argument}
The section variables are the instance \texttt{[Finite\allowbreak{}Dimensional \ensuremath{\mathbb R} E]} and the hypotheses \texttt{hf : Convex\allowbreak{}On \ensuremath{\mathbb R} univ f}, \texttt{hb : \ensuremath{\exists} b, \ensuremath{\forall} z, b \ensuremath{\le} f z} and \texttt{hlam : 0 < lam}. The command \texttt{include hf hb hlam} makes these three hypotheses arguments of every theorem of the section, and a line \texttt{omit … in} removes some of them from the single theorem that follows it.




\texttt{moreau f lam x} is defined as \texttt{prox\allowbreak{}Objective f lam x (prox f lam x)}, the objective evaluated at the selected proximal point, and \texttt{residual f lam x} is \texttt{lam\ensuremath{{}^{-1}} \ensuremath{\bullet} (x -\allowbreak{} prox f lam x)}. Under the section hypotheses \texttt{prox f lam x} is the unique proximal point \(P_\lambda(x)\) (by \texttt{prox\_\allowbreak{}is\allowbreak{}Prox\allowbreak{}Point}), so these definitions are \(f^\lambda\) and \(r_\lambda\).




\texttt{moreau\_\allowbreak{}le} is (a): it is \texttt{prox\_\allowbreak{}is\allowbreak{}Prox\allowbreak{}Point hf hb hlam x} tested at the competitor \texttt{x}, simplified by \texttt{simpa} since \(\|x-x\|^2/(2\lambda)=0\). \texttt{le\_\allowbreak{}moreau} is (b). Its line \texttt{omit [Inner\allowbreak{}Product\allowbreak{}Space \ensuremath{\mathbb R} E] [Finite\allowbreak{}Dimensional \ensuremath{\mathbb R} E] hf hb in} removes the inner-product structure, finite dimension, convexity and \texttt{hb}, so the theorem holds for every function with the lower bound \texttt{hbf : \ensuremath{\forall} z, b \ensuremath{\le} f z} and every \(\lambda>0\). The proof chains \texttt{hbf \_\allowbreak{}} with \texttt{le\_\allowbreak{}add\_\allowbreak{}of\_\allowbreak{}nonneg\_\allowbreak{}right}, the nonnegativity of the penalty being proved by \texttt{positivity}.
\end{leanbox}

\section{\texorpdfstring{Global supports, the quadratic upper bound and convexity}{Global supports, the quadratic upper bound and convexity}}


\textbf{Claim.} For all \(x,y\in E\),


\[f^\lambda(x)+\langle r_\lambda(x),y-x\rangle\le f^\lambda(y)\le f^\lambda(x)+\langle r_\lambda(x),y-x\rangle+\frac{\|y-x\|^2}{2\lambda}.\qquad(1)\]


In particular \(r_\lambda(x)\) is a subgradient of \(f^\lambda\) at \(x\), and \(f^\lambda\) is convex.




\textbf{Proof of the lower bound.} Write \(p=P_\lambda(x)\) and \(q=P_\lambda(y)\). By the residual subgradient condition of the previous section (here convexity of \(f\) is used), \(\lambda^{-1}(x-p)\in\partial f(p)\). Testing this support at \(q\),


\[f(q)\ge f(p)+\lambda^{-1}\langle x-p,q-p\rangle.\qquad(2)\]


Put \(a=q-y\) and \(c=p-x\). Then \(q-p=a-c+(y-x)\) and \(x-p=-c\), so


\[\langle x-p,q-p\rangle=-\langle c,a\rangle+\|c\|^2+\langle x-p,y-x\rangle.\]


Insert this into (2), add \(\|a\|^2/(2\lambda)\) to both sides and subtract \(\|c\|^2/(2\lambda)+\lambda^{-1}\langle x-p,y-x\rangle\):


\[f(q)+\frac{\|a\|^2}{2\lambda}-\Bigl(f(p)+\frac{\|c\|^2}{2\lambda}+\frac1\lambda\langle x-p,y-x\rangle\Bigr)\ge\frac{\|a\|^2-2\langle a,c\rangle+\|c\|^2}{2\lambda}=\frac{\|a-c\|^2}{2\lambda}\ge0.\]


Since \(f(q)+\|a\|^2/(2\lambda)=f^\lambda(y)\), \(f(p)+\|c\|^2/(2\lambda)=f^\lambda(x)\) and \(\lambda^{-1}\langle x-p,y-x\rangle=\langle r_\lambda(x),y-x\rangle\), this is the lower bound in (1).




\textbf{Proof of the upper bound.} Since \(P_\lambda(y)\) minimizes \(Q_{\lambda,y}\), the competitor \(p=P_\lambda(x)\) gives \(f^\lambda(y)\le Q_{\lambda,y}(p)=f(p)+\|p-y\|^2/(2\lambda)\). Writing \(p-y=(p-x)-(y-x)\),


\[\|p-y\|^2=\|p-x\|^2-2\langle p-x,y-x\rangle+\|y-x\|^2,\]


and therefore


\[Q_{\lambda,y}(p)=f(p)+\frac{\|p-x\|^2}{2\lambda}+\frac1\lambda\langle x-p,y-x\rangle+\frac{\|y-x\|^2}{2\lambda}=f^\lambda(x)+\langle r_\lambda(x),y-x\rangle+\frac{\|y-x\|^2}{2\lambda}.\]


This is the upper bound in (1); it uses only the minimality of \(P_\lambda(y)\).




\textbf{Proof of convexity.} Let \(x,y\in E\), \(a,b\ge0\) with \(a+b=1\), and \(z=ax+by\). The lower bound in (1), based at \(z\) and evaluated at \(x\) and at \(y\), gives


\[f^\lambda(z)+\langle r_\lambda(z),x-z\rangle\le f^\lambda(x),\qquad f^\lambda(z)+\langle r_\lambda(z),y-z\rangle\le f^\lambda(y).\]


Multiply the first inequality by \(a\ge0\), the second by \(b\ge0\), and add. The inner-product terms combine to \(\langle r_\lambda(z),a(x-z)+b(y-z)\rangle\), and \(a(x-z)+b(y-z)=ax+by-(a+b)z=z-z=0\). Since \((a+b)f^\lambda(z)=f^\lambda(z)\), the sum reads \(f^\lambda(ax+by)\le af^\lambda(x)+bf^\lambda(y)\). \(\square\)

\begin{leanbox}{Lean: \texttt{moreau\_\allowbreak{}support}, \texttt{moreau\_\allowbreak{}upper}, \texttt{moreau\_\allowbreak{}convex}}
\leansource{\texttt{moreau\_\allowbreak{}support} \quad ProximalLangevin/\allowbreak{}Moreau.lean, lines 47–66}
\begin{leancode}
\leanline{0}{theorem\ moreau\_support\ (x\ y\ :\ E)\ :}
\leanline{4}{moreau\ f\ lam\ x\ +\ {\leanfallback ⟪}residual\ f\ lam\ x,\ y\ -\ x{\leanfallback ⟫}\ ≤\ moreau\ f\ lam\ y}
\end{leancode}
\leansource{\texttt{moreau\_\allowbreak{}upper} \quad ProximalLangevin/\allowbreak{}Moreau.lean, lines 69–79}
\begin{leancode}
\leanline{0}{theorem\ moreau\_upper\ (x\ y\ :\ E)\ :}
\leanline{4}{moreau\ f\ lam\ y\ ≤\ moreau\ f\ lam\ x\ +\ {\leanfallback ⟪}residual\ f\ lam\ x,\ y\ -\ x{\leanfallback ⟫}\ +}
\leanline{6}{‖y\ -\ x‖\ \textasciicircum{}\ 2\ /\ (2\ *\ lam)}
\end{leancode}
\leansource{\texttt{moreau\_\allowbreak{}convex} \quad ProximalLangevin/\allowbreak{}Moreau.lean, lines 82–103}
\begin{leancode}
\leanline{0}{theorem\ moreau\_convex\ :\ ConvexOn\ ℝ\ univ\ (moreau\ f\ lam)}
\end{leancode}
\leanlabel{How the Lean proof follows the argument}
\texttt{moreau\_\allowbreak{}support} is the lower bound in (1). With \texttt{p := prox f lam x} and \texttt{q := prox f lam y}, \texttt{h} is (2): \texttt{subgradient\_\allowbreak{}of\_\allowbreak{}is\allowbreak{}Prox\allowbreak{}Point} applied to \texttt{prox\_\allowbreak{}is\allowbreak{}Prox\allowbreak{}Point hf hb hlam x} and tested at \texttt{q}. \texttt{hd} is the decomposition \(q-p=(q-y)-(p-x)+(y-x)\), and the rewrites \texttt{inner\_\allowbreak{}add\_\allowbreak{}right}, \texttt{inner\_\allowbreak{}sub\_\allowbreak{}right}, \texttt{inner\_\allowbreak{}neg\_\allowbreak{}left} and \texttt{real\_\allowbreak{}inner\_\allowbreak{}self\_\allowbreak{}eq\_\allowbreak{}norm\_\allowbreak{}sq} expand the inner product as in the proof. \texttt{hs} is \(\|a-c\|^2\ge0\), expanded by \texttt{norm\_\allowbreak{}sub\_\allowbreak{}sq\_\allowbreak{}real} and \texttt{real\_\allowbreak{}inner\_\allowbreak{}comm}. The goal, displayed by \texttt{change}, is multiplied by \(2\lambda>0\) using \texttt{mul\_\allowbreak{}le\_\allowbreak{}mul\_\allowbreak{}iff\_\allowbreak{}of\_\allowbreak{}pos\_\allowbreak{}right}; \texttt{h'} is (2) multiplied by \(\lambda\), and \texttt{field\_\allowbreak{}simp} and \texttt{nlinarith} combine \texttt{h'} with \texttt{hs}.




\texttt{moreau\_\allowbreak{}upper} is the upper bound in (1). \texttt{h} is the minimality of \texttt{prox f lam y} tested at the competitor \texttt{prox f lam x}; \texttt{h.\allowbreak{}trans\_\allowbreak{}eq} reduces the goal to the displayed identity for \(Q_{\lambda,y}(p)\), which follows from the rewrite \(p-y=(p-x)-(y-x)\), \texttt{norm\_\allowbreak{}sub\_\allowbreak{}sq\_\allowbreak{}real}, \texttt{real\_\allowbreak{}inner\_\allowbreak{}smul\_\allowbreak{}left}, a sign change, \texttt{field\_\allowbreak{}simp} and \texttt{ring}.




\texttt{moreau\_\allowbreak{}convex} proves \texttt{Convex\allowbreak{}On \ensuremath{\mathbb R} univ (moreau f lam)}. This structure has two components: \texttt{convex\_\allowbreak{}univ} (the domain is convex) and the convexity inequality for weights \texttt{a}, \texttt{b} with \texttt{ha : 0 \ensuremath{\le} a}, \texttt{hb' : 0 \ensuremath{\le} b} and \texttt{hab : a + b = 1}. \texttt{hx} and \texttt{hy} are the two supports based at \texttt{z}; \texttt{hi} states that the weighted inner-product terms sum to zero, using \texttt{hz}, the identity \(a(x-z)+b(y-z)=0\), and \texttt{inner\_\allowbreak{}zero\_\allowbreak{}right}. \texttt{h1} and \texttt{h2} are the weighted supports, \texttt{hm} is \(af^\lambda(z)+bf^\lambda(z)=f^\lambda(z)\), and \texttt{nlinarith} adds them.
\end{leanbox}

\section{\texorpdfstring{Lipschitz bounds for the proximal map and the residual}{Lipschitz bounds for the proximal map and the residual}}


\textbf{Claim.} For all \(x,y\in E\),


\[\|P_\lambda(x)-P_\lambda(y)\|\le\|x-y\|,\qquad\|r_\lambda(x)-r_\lambda(y)\|\le\frac1\lambda\|x-y\|.\]


\textbf{Proof.} The first inequality is the nonexpansiveness of proximal points from the previous section, applied to \(p=P_\lambda(x)\) and \(q=P_\lambda(y)\). For the second, \(r_\lambda(x)-r_\lambda(y)=\lambda^{-1}\bigl((x-P_\lambda(x))-(y-P_\lambda(y))\bigr)\). By homogeneity of the norm and \(\lambda^{-1}>0\),


\[\|r_\lambda(x)-r_\lambda(y)\|=\lambda^{-1}\bigl\|(x-P_\lambda(x))-(y-P_\lambda(y))\bigr\|\le\lambda^{-1}\|x-y\|,\]


where the inequality is the nonexpansiveness of the residual map \(x\mapsto x-P_\lambda(x)\) from the previous section. \(\square\)

\begin{leanbox}{Lean: \texttt{prox\_\allowbreak{}lipschitz}, \texttt{residual\_\allowbreak{}lipschitz}}
\leansource{\texttt{prox\_\allowbreak{}lipschitz} \quad ProximalLangevin/\allowbreak{}Moreau.lean, lines 106–110}
\begin{leancode}
\leanline{0}{theorem\ prox\_lipschitz\ :\ LipschitzWith\ 1\ (prox\ f\ lam)}
\end{leancode}
\leansource{\texttt{residual\_\allowbreak{}lipschitz} \quad ProximalLangevin/\allowbreak{}Moreau.lean, lines 113–120}
\begin{leancode}
\leanline{0}{theorem\ residual\_lipschitz\ :\ LipschitzWith\ lam⁻¹.toNNReal\ (residual\ f\ lam)}
\end{leancode}
\leanlabel{How the Lean proof follows the argument}
\texttt{Lipschitz\allowbreak{}With K g} means \(\operatorname{dist}(g(x),g(y))\le K\operatorname{dist}(x,y)\) for all \(x,y\), with \(K\) a nonnegative real number (type \texttt{NNReal}, also written \texttt{\ensuremath{\mathbb R}\ensuremath{\ge}0}); in a normed space \(\operatorname{dist}(u,v)=\|u-v\|\) (\texttt{dist\_\allowbreak{}eq\_\allowbreak{}norm\_\allowbreak{}sub}).




\texttt{prox\_\allowbreak{}lipschitz} states \texttt{Lipschitz\allowbreak{}With 1 (prox f lam)}. After \texttt{Lipschitz\allowbreak{}With.\allowbreak{}mk\_\allowbreak{}one} the goal is the inequality between distances, which is \texttt{Is\allowbreak{}Prox\allowbreak{}Point.\allowbreak{}nonexpansive} applied to the two proximal points given by \texttt{prox\_\allowbreak{}is\allowbreak{}Prox\allowbreak{}Point}. \texttt{residual\_\allowbreak{}lipschitz} states \texttt{Lipschitz\allowbreak{}With lam\ensuremath{{}^{-1}}.\allowbreak{}to\allowbreak{}NNReal (residual f lam)}, where \texttt{lam\ensuremath{{}^{-1}}.\allowbreak{}to\allowbreak{}NNReal} is \(\lambda^{-1}\) regarded as a nonnegative real. \texttt{Lipschitz\allowbreak{}With.\allowbreak{}of\_\allowbreak{}dist\_\allowbreak{}le'} reduces it to \(\|r_\lambda(x)-r_\lambda(y)\|\le\lambda^{-1}\|x-y\|\); the rewrites \texttt{← smul\_\allowbreak{}sub}, \texttt{norm\_\allowbreak{}smul}, \texttt{Real.\allowbreak{}norm\_\allowbreak{}eq\_\allowbreak{}abs} and \texttt{abs\_\allowbreak{}of\_\allowbreak{}pos (inv\_\allowbreak{}pos.\allowbreak{}mpr hlam)} extract the factor \(\lambda^{-1}\), and \texttt{mul\_\allowbreak{}le\_\allowbreak{}mul\_\allowbreak{}of\_\allowbreak{}nonneg\_\allowbreak{}left} multiplies \texttt{Is\allowbreak{}Prox\allowbreak{}Point.\allowbreak{}residual\_\allowbreak{}nonexpansive} by it.
\end{leanbox}

\section{\texorpdfstring{The gradient of the envelope is the scaled residual}{The gradient of the envelope is the scaled residual}}


Recall that a function \(g:E\to\mathbb R\) has gradient \(v\in E\) at \(x\) if \(g(y)-g(x)-\langle v,y-x\rangle=o(\|y-x\|)\) as \(y\to x\): for every \(c>0\) there is \(\rho>0\) such that


\[|g(y)-g(x)-\langle v,y-x\rangle|\le c\|y-x\|\qquad\text{whenever }\|y-x\|<\rho.\]


Equivalently, \(g\) is Fréchet differentiable at \(x\) with derivative \(h\mapsto\langle v,h\rangle\); \(v\) is then denoted \(\nabla g(x)\).




\textbf{Claim.} For every \(x\in E\), \(f^\lambda\) has gradient \(r_\lambda(x)=\lambda^{-1}(x-P_\lambda(x))\) at \(x\). Consequently \(f^\lambda\) is differentiable on \(E\), \(\nabla f^\lambda=r_\lambda\), and \(f^\lambda\) is continuously differentiable. No differentiability of \(f\) is assumed.




\textbf{Proof.} Fix \(x\) and put \(R(y)=f^\lambda(y)-f^\lambda(x)-\langle r_\lambda(x),y-x\rangle\). By (1) of the previous step but one,


\[0\le R(y)\le\frac{\|y-x\|^2}{2\lambda}\qquad(y\in E).\qquad(1)\]


Let \(c>0\) and \(\rho=2\lambda c>0\). If \(\|y-x\|<\rho\), then by (1)


\[|R(y)|=R(y)\le\frac{\|y-x\|}{2\lambda}\,\|y-x\|\le\frac{\rho}{2\lambda}\,\|y-x\|=c\,\|y-x\|.\]


Hence \(f^\lambda\) has gradient \(r_\lambda(x)\) at \(x\); in particular it is differentiable at every point and \(\nabla f^\lambda(x)=r_\lambda(x)\).




For the continuity of the derivative, the derivative at \(x\) is the linear functional \(h\mapsto\langle r_\lambda(x),h\rangle\). The map \(v\mapsto\langle v,\cdot\rangle\) from \(E\) to its dual is a linear isometry (Riesz representation), hence continuous, and \(r_\lambda\) is \(\lambda^{-1}\)-Lipschitz by the previous step, hence continuous. Their composition \(x\mapsto\langle r_\lambda(x),\cdot\rangle\) is continuous. A function that is differentiable everywhere with continuous derivative is continuously differentiable. \(\square\)




By the previous step the gradient \(\nabla f^\lambda\) is \(\lambda^{-1}\)-Lipschitz.

\begin{leanbox}{Lean: \texttt{moreau\_\allowbreak{}hasGradientAt}, \texttt{moreau\_\allowbreak{}differentiable}, \texttt{gradient\_\allowbreak{}moreau} and 1 more}
\leansource{\texttt{moreau\_\allowbreak{}hasGradientAt} \quad ProximalLangevin/\allowbreak{}Moreau.lean, lines 123–136}
\begin{leancode}
\leanline{0}{theorem\ moreau\_hasGradientAt\ (x\ :\ E)\ :\ HasGradientAt\ (moreau\ f\ lam)\ (residual\ f\ lam\ x)\ x}
\end{leancode}
\leansource{\texttt{moreau\_\allowbreak{}differentiable} \quad ProximalLangevin/\allowbreak{}Moreau.lean, lines 139–140}
\begin{leancode}
\leanline{0}{theorem\ moreau\_differentiable\ :\ Differentiable\ ℝ\ (moreau\ f\ lam)}
\end{leancode}
\leansource{\texttt{gradient\_\allowbreak{}moreau} \quad ProximalLangevin/\allowbreak{}Moreau.lean, lines 143–144}
\begin{leancode}
\leanline{0}{theorem\ gradient\_moreau\ (x\ :\ E)\ :\ gradient\ (moreau\ f\ lam)\ x\ =\ residual\ f\ lam\ x}
\end{leancode}
\leansource{\texttt{moreau\_\allowbreak{}contDiff} \quad ProximalLangevin/\allowbreak{}Moreau.lean, lines 147–153}
\begin{leancode}
\leanline{0}{theorem\ moreau\_contDiff\ :\ ContDiff\ ℝ\ 1\ (moreau\ f\ lam)}
\end{leancode}
\leanlabel{How the Lean proof follows the argument}
\texttt{Has\allowbreak{}Gradient\allowbreak{}At g v x} states that \(g\) is Fréchet differentiable at \(x\) with derivative \(h\mapsto\langle v,h\rangle\), and \texttt{gradient g x} is the vector representing the derivative of \(g\) at \(x\).



\begin{enumerate}


\item \texttt{moreau\_\allowbreak{}has\allowbreak{}Gradient\allowbreak{}At}: \texttt{has\allowbreak{}Gradient\allowbreak{}At\_\allowbreak{}iff\_\allowbreak{}is\allowbreak{}Little\allowbreak{}O} rewrites the goal as the statement that \(R(y)\) is little-o of \(y-x\) along the neighbourhood filter \texttt{{\leanfallback 𝓝} x}, and \texttt{Asymptotics.\allowbreak{}is\allowbreak{}Little\allowbreak{}O\_\allowbreak{}iff} turns this into: for every \(c>0\), eventually \(|R(y)|\le c\|y-x\|\). \texttt{filter\_\allowbreak{}upwards [Metric.\allowbreak{}ball\_\allowbreak{}mem\_\allowbreak{}nhds x …]} restricts to the ball of radius \(c\cdot(2\lambda)\), which gives \texttt{hdist}. \texttt{hlo} (\texttt{moreau\_\allowbreak{}support}) and \texttt{hhi} (\texttt{moreau\_\allowbreak{}upper}) are the two sides of (1); \texttt{abs\_\allowbreak{}of\_\allowbreak{}nonneg} removes the absolute value using the lower side, and \texttt{div\_\allowbreak{}le\_\allowbreak{}iff\ensuremath{{}_{0}}} with \texttt{nlinarith} performs the final estimate.

\item \texttt{moreau\_\allowbreak{}differentiable} applies \texttt{.\allowbreak{}differentiable\allowbreak{}At} to the gradient statement at each point, and \texttt{gradient\_\allowbreak{}moreau}, the identity \(\nabla f^\lambda(x)=r_\lambda(x)\), is \texttt{Has\allowbreak{}Gradient\allowbreak{}At.\allowbreak{}gradient} applied to it.

\item \texttt{moreau\_\allowbreak{}cont\allowbreak{}Diff} proves \texttt{Cont\allowbreak{}Diff \ensuremath{\mathbb R} 1 (moreau f lam)}. By \texttt{cont\allowbreak{}Diff\_\allowbreak{}one\_\allowbreak{}iff\_\allowbreak{}has\allowbreak{}FDeriv\allowbreak{}At} it suffices to give a continuous family of Fréchet derivatives. The family is \texttt{Inner\allowbreak{}Product\allowbreak{}Space.\allowbreak{}to\allowbreak{}Dual \ensuremath{\mathbb R} E (residual f lam x)}, the functional \(\langle r_\lambda(x),\cdot\rangle\); it is continuous as the composition of the continuous map \texttt{Inner\allowbreak{}Product\allowbreak{}Space.\allowbreak{}to\allowbreak{}Dual \ensuremath{\mathbb R} E} with the residual, which is continuous by \texttt{residual\_\allowbreak{}lipschitz}, and each derivative statement is \texttt{(moreau\_\allowbreak{}has\allowbreak{}Gradient\allowbreak{}At hf hb hlam x).\allowbreak{}has\allowbreak{}FDeriv\allowbreak{}At}.

\end{enumerate}
\end{leanbox}

\section{\texorpdfstring{Bounded gradient and bounded proximal displacement}{Bounded gradient and bounded proximal displacement}}


\textbf{Claim.} Suppose in addition that \(f\) is \(K\)-Lipschitz for some \(K\ge0\), that is, \(|f(u)-f(v)|\le K\|u-v\|\) for all \(u,v\in E\). Then for every \(x\in E\),


\[\|r_\lambda(x)\|\le K,\qquad\|x-P_\lambda(x)\|\le\lambda K.\]


\textbf{Proof.} Put \(p=P_\lambda(x)\) and \(g=r_\lambda(x)\). By the residual subgradient condition, \(g\in\partial f(p)\). Testing the support inequality at \(y=p+g\) gives \(f(p)+\langle g,g\rangle\le f(p+g)\), so, by the Lipschitz property,


\[\|g\|^2\le f(p+g)-f(p)\le|f(p+g)-f(p)|\le K\|g\|.\]


If \(\|g\|>K\), then \(\|g\|>0\) and \(\|g\|^2>K\|g\|\), contradicting the display; hence \(\|g\|\le K\). Since \(x-P_\lambda(x)=\lambda g\) and \(\lambda>0\), \(\|x-P_\lambda(x)\|=\lambda\|g\|\le\lambda K\). \(\square\)




Because \(\nabla f^\lambda=r_\lambda\), the first inequality bounds the gradient of the envelope by \(K\), uniformly in \(\lambda>0\).

\begin{leanbox}{Lean: \texttt{residual\_\allowbreak{}norm\_\allowbreak{}le}, \texttt{prox\_\allowbreak{}displacement\_\allowbreak{}le}}
\leansource{\texttt{residual\_\allowbreak{}norm\_\allowbreak{}le} \quad ProximalLangevin/\allowbreak{}Moreau.lean, lines 156–167}
\begin{leancode}
\leanline{0}{theorem\ residual\_norm\_le\ \{K\ :\ NNReal\}\ (hK\ :\ LipschitzWith\ K\ f)\ (x\ :\ E)\ :}
\leanline{4}{‖residual\ f\ lam\ x‖\ ≤\ (K\ :\ ℝ)}
\end{leancode}
\leansource{\texttt{prox\_\allowbreak{}displacement\_\allowbreak{}le} \quad ProximalLangevin/\allowbreak{}Moreau.lean, lines 170–176}
\begin{leancode}
\leanline{0}{theorem\ prox\_displacement\_le\ \{K\ :\ NNReal\}\ (hK\ :\ LipschitzWith\ K\ f)\ (x\ :\ E)\ :}
\leanline{4}{‖x\ -\ prox\ f\ lam\ x‖\ ≤\ lam\ *\ (K\ :\ ℝ)}
\end{leancode}
\leanlabel{How the Lean proof follows the argument}
The Lipschitz hypothesis is \texttt{h\allowbreak{}K : Lipschitz\allowbreak{}With K f} with \texttt{K : NNReal}, and \texttt{(K : \ensuremath{\mathbb R})} is the coercion of \(K\) to a real number. In \texttt{residual\_\allowbreak{}norm\_\allowbreak{}le}, \texttt{hs} is the support from \texttt{subgradient\_\allowbreak{}of\_\allowbreak{}is\allowbreak{}Prox\allowbreak{}Point} tested at \texttt{p + g}; \texttt{add\_\allowbreak{}sub\_\allowbreak{}cancel\_\allowbreak{}left} simplifies \((p+g)-p\) to \(g\) and \texttt{real\_\allowbreak{}inner\_\allowbreak{}self\_\allowbreak{}eq\_\allowbreak{}norm\_\allowbreak{}sq} rewrites \(\langle g,g\rangle\) as \(\|g\|^2\). \texttt{hd} is the Lipschitz bound \texttt{h\allowbreak{}K.\allowbreak{}dist\_\allowbreak{}le\_\allowbreak{}mul (p + g) p}, rewritten with \texttt{Real.\allowbreak{}dist\_\allowbreak{}eq} (on \(\mathbb R\), \(\operatorname{dist}(a,b)=|a-b|\)), and \texttt{hdiff} drops the absolute value using \texttt{le\_\allowbreak{}abs\_\allowbreak{}self}. \texttt{nlinarith}, given \(\|g\|\ge0\) and \(K\ge0\), carries out the case argument.




\texttt{prox\_\allowbreak{}displacement\_\allowbreak{}le} rewrites the residual bound with \texttt{norm\_\allowbreak{}smul}, \texttt{Real.\allowbreak{}norm\_\allowbreak{}eq\_\allowbreak{}abs} and \texttt{abs\_\allowbreak{}of\_\allowbreak{}pos}, giving \(\lambda^{-1}\|x-P_\lambda(x)\|\le K\); it multiplies this by \(\lambda\ge0\) (\texttt{mul\_\allowbreak{}le\_\allowbreak{}mul\_\allowbreak{}of\_\allowbreak{}nonneg\_\allowbreak{}left}), and \texttt{field\_\allowbreak{}simp} cancels \(\lambda\lambda^{-1}\).
\end{leanbox}

\section{\texorpdfstring{The uniform gap between a Lipschitz function and its envelope}{The uniform gap between a Lipschitz function and its envelope}}


\textbf{Claim.} (a) Let \(f:E\to\mathbb R\) be \(K\)-Lipschitz with \(K\ge0\), and \(\lambda>0\). Then for all \(x,p\in E\)


\[Q_{\lambda,x}(p)=f(p)+\frac{\|p-x\|^2}{2\lambda}\ge f(x)-\frac{\lambda K^2}{2};\]


in particular \(f^\lambda(x)\ge f(x)-\lambda K^2/2\). Part (a) uses neither convexity nor a lower bound. (b) If \(f\) is in addition convex and bounded below, then for every \(x\in E\)


\[0\le f(x)-f^\lambda(x)\le\frac{\lambda K^2}{2}.\]


\textbf{Proof.} (a) Let \(d=\|p-x\|\). The Lipschitz property gives \(f(x)-f(p)\le|f(x)-f(p)|\le Kd\), i.e. \(f(p)\ge f(x)-Kd\). Completing the square in \(d\),


\[\frac{d^2}{2\lambda}-Kd=\frac{(d-\lambda K)^2}{2\lambda}-\frac{\lambda K^2}{2}\ge-\frac{\lambda K^2}{2},\]


because \((d-\lambda K)^2\ge0\) and \(2\lambda>0\). Hence \(Q_{\lambda,x}(p)\ge f(x)-Kd+d^2/(2\lambda)\ge f(x)-\lambda K^2/2\).




(b) The inequality \(f(x)-f^\lambda(x)\ge0\) is part (a) of the first step, and \(f(x)-f^\lambda(x)\le\lambda K^2/2\) is part (a) of this step with \(p=P_\lambda(x)\). \(\square\)




The gap is uniform in \(x\) and linear in \(\lambda\). For any function \(g_1\), writing \(U=g_1+f\) and \(V_\lambda=g_1+f^\lambda\), it gives \(V_\lambda\le U\le V_\lambda+\lambda K^2/2\) on all of \(E\); this two-sided comparison is the form in which the gap enters the comparison of marginal potentials.

\begin{leanbox}{Lean: \texttt{moreau\_\allowbreak{}lower\_\allowbreak{}of\_\allowbreak{}lipschitz}, \texttt{moreau\_\allowbreak{}gap\_\allowbreak{}of\_\allowbreak{}lipschitz}}
\leansource{\texttt{moreau\_\allowbreak{}lower\_\allowbreak{}of\_\allowbreak{}lipschitz} \quad ProximalLangevin/\allowbreak{}Moreau.lean, lines 180–190}
\begin{leancode}
\leanline{0}{theorem\ moreau\_lower\_of\_lipschitz\ \{K\ :\ NNReal\}\ (hK\ :\ LipschitzWith\ K\ f)\ (x\ :\ E)\ :}
\leanline{4}{f\ x\ -\ lam\ *\ (K\ :\ ℝ)\ \textasciicircum{}\ 2\ /\ 2\ ≤\ moreau\ f\ lam\ x}
\end{leancode}
\leansource{\texttt{moreau\_\allowbreak{}gap\_\allowbreak{}of\_\allowbreak{}lipschitz} \quad ProximalLangevin/\allowbreak{}Moreau.lean, lines 193–200}
\begin{leancode}
\leanline{0}{theorem\ moreau\_gap\_of\_lipschitz\ \{K\ :\ NNReal\}\ (hK\ :\ LipschitzWith\ K\ f)\ (x\ :\ E)\ :}
\leanline{4}{0\ ≤\ f\ x\ -\ moreau\ f\ lam\ x\ ∧\ f\ x\ -\ moreau\ f\ lam\ x\ ≤\ lam\ *\ (K\ :\ ℝ)\ \textasciicircum{}\ 2\ /\ 2}
\end{leancode}
\leanlabel{How the Lean proof follows the argument}
\texttt{moreau\_\allowbreak{}lower\_\allowbreak{}of\_\allowbreak{}lipschitz} is preceded by \texttt{omit [Inner\allowbreak{}Product\allowbreak{}Space \ensuremath{\mathbb R} E] [Finite\allowbreak{}Dimensional \ensuremath{\mathbb R} E] hf hb in}, so its only hypotheses are \texttt{hlam} and \texttt{h\allowbreak{}K}. It therefore holds for every \(K\)-Lipschitz function and for whichever point \texttt{prox f lam x} denotes, which is the content of (a). \texttt{hdist} is the Lipschitz bound between \texttt{x} and \texttt{prox f lam x}, with \texttt{norm\_\allowbreak{}sub\_\allowbreak{}rev} orienting the norm, and \texttt{hdiff} removes the absolute value. After multiplying the goal by \(2\lambda>0\) and \texttt{field\_\allowbreak{}simp}, \texttt{nlinarith} uses \texttt{hdiff'} (\texttt{hdiff} multiplied by \(\lambda\)) and the square \((d-\lambda K)^2\ge0\), supplied as \texttt{sq\_\allowbreak{}nonneg (\ensuremath{\Vert}prox f lam x -\allowbreak{} x\ensuremath{\Vert} -\allowbreak{} lam * (K : \ensuremath{\mathbb R}))}: this is the completion of the square.




\texttt{moreau\_\allowbreak{}gap\_\allowbreak{}of\_\allowbreak{}lipschitz} is (b), a conjunction \texttt{∧} proved by \texttt{constructor}. The first half is \texttt{moreau\_\allowbreak{}le} turned into \(0\le f(x)-f^\lambda(x)\) by \texttt{sub\_\allowbreak{}nonneg.\allowbreak{}mpr}; the second half is \texttt{moreau\_\allowbreak{}lower\_\allowbreak{}of\_\allowbreak{}lipschitz} rearranged by \texttt{linarith}.
\end{leanbox}

\end{document}
